% TOP_PROBLEM: {"id":30007641,"problem_number":"LOCAL-30007641","title":"Temporary contracts: stepping stones or traps","category_id":21,"category":{"id":21,"name":"economics","display_name":"Economics","description":"Open economic theory statements, published research agendas, and proposed scoped specifications; question provenance and historical priority are qualified per record.","slug":"economics","order_index":21,"created_at":"2026-10-08T00:00:00Z"},"status":"research_question","external_url":null,"legacy_ids":[],"published":false,"statement":"In the explicitly specified setting: Which employment-protection reforms make temporary jobs reliable routes to stable careers rather than persistent cycles of insecurity and weak training? The mathematical statement above fixes the target, assumptions and scope of this catalogue formulation.","economics":{"original_id":130,"field":"Labor markets and work","kind":"identification","formalization_basis":"adapted_research_question","source_status":"research_agenda","priority_status":"earliest_found","priority_note":"Earliest directly inspected qualifying passage in this audit; earlier origins were not established. A review or research-agenda statement does not imply original priority.","earliest_verified_reference":{"authors":["Alison L. Booth","Marco Francesconi","Jeff Frank"],"title":"Temporary Jobs: Stepping Stones or Dead Ends?","year":2000,"url":"https://docs.iza.org/dp205.pdf","doi":"10.1111/1468-0297.00043","locator":"IZA DP 205, section I, Introduction, printed pp. 1–2","evidence_summary":"The introduction asks whether temporary jobs provide good careers or routes to permanent work, and then answers for 1990s Britain."},"current_reference":{"authors":["Mattia Filomena","Matteo Picchio"],"title":"Are Temporary Jobs Stepping Stones or Dead Ends? A Meta-Analytical Review of the Literature","year":2021,"url":"https://docs.iza.org/dp14367.pdf","doi":"10.1108/IJM-02-2022-0064","locator":"Section 1, Introduction, printed pp. 1–2; section 2, printed pp. 2–3","evidence_summary":"The review documents conflicting causal findings and context-dependent theoretical predictions, without a universal sign conclusion."},"formalization_note":"The cited passage establishes a research direction, not this exact mathematical statement. The notation, population, policy menu, assumptions, and estimands are an original adaptation for this catalogue; no universal unsolved theorem or source priority is claimed."},"statusQualification":"Published question or research agenda verified at the cited locator. The mathematical specification is adapted for this catalogue and includes additional assumptions; source publication does not establish first-ever priority or certify this exact model remains unsolved. The cited passage establishes a research direction, not this exact mathematical statement. The notation, population, policy menu, assumptions, and estimands are an original adaptation for this catalogue; no universal unsolved theorem or source priority is claimed.","statusReviewedAt":"2026-10-08"}
% FIRST_POSTED: 2026-10-08
% ENTRY_KIND: statement_only
% !TeX program = lualatex
\documentclass[11pt]{article}
\usepackage{fontspec}
\usepackage[a4paper,margin=25mm]{geometry}
\usepackage{amsmath,amssymb,mathtools}
\usepackage{xurl}
\usepackage[hidelinks]{hyperref}
\newcommand{\sref}[2]{\hyperlink{source:#1:#2}{[S#2]}}
\setlength{\emergencystretch}{3em}
\begin{document}
% problemId:30007641
\section*{Temporary contracts: stepping stones or traps}\label{30007641}
\subsection{Definitions and mathematical statement}
Fix a cohort of jobseekers, a labor market, and an employment horizon \(T\). Let \(D_i=1\) denote acceptance of a precisely specified temporary contract and \(D_i=0\) the declared alternative, such as continued job search rather than immediate permanent employment. Let \(A_{it}(d)\) indicate stable employment and \(Y_{it}(d)\) real earnings at date \(t\). Define \(\tau_A(t,x)=E[A_{it}(1)-A_{it}(0)\mid X_i=x]\) and \(\tau_Y(T,x)=E[\sum_{t=0}^{T}\delta^t\{Y_{it}(1)-Y_{it}(0)\}\mid X_i=x]\), with baseline characteristics \(X_i\) and discount factor \(\delta\in(0,1]\). A stepping stone requires a favorable declared career outcome relative to this comparator, not merely an observed transition to a permanent job. Identify these effects using credible assignment or contract-offer variation, allowing selection and duration dependence. Determine how results vary by contract type, worker disadvantage, demand conditions, and employment-protection rules. Evaluate reforms separately as market-wide policies because firms can substitute among contracts. Earlier studies answer specific settings; the research agenda concerns reliable context-dependent effects and mechanisms, not a conjecture that every temporary job has one common sign.

\par\textbf{Formulation and scope.} The cited passage establishes a research direction, not this exact mathematical statement. The notation, population, policy menu, assumptions, and estimands are an original adaptation for this catalogue; no universal unsolved theorem or source priority is claimed.

\par\textbf{Open-question provenance.} An explicit question or research agenda was verified at the source locator. This does not by itself certify that every later version remains unresolved \sref{30007641}{1}.

\par\textbf{Historical priority.} Earliest directly inspected qualifying passage in this audit; earlier origins were not established. A review or research-agenda statement does not imply original priority.

\subsection{Short English statement}
In the explicitly specified setting: Which employment-protection reforms make temporary jobs reliable routes to stable careers rather than persistent cycles of insecurity and weak training? The mathematical statement above fixes the target, assumptions and scope of this catalogue formulation.

\subsection{Sources}
\begin{itemize}\raggedright
\item[\textnormal{[S1]}]\hypertarget{source:30007641:1}{} Alison L. Booth, Marco Francesconi, Jeff Frank. 2000. Temporary Jobs: Stepping Stones or Dead Ends?. IZA DP 205, section I, Introduction, printed pp. 1–2.
\url{https://docs.iza.org/dp205.pdf}.
DOI: 10.1111/1468-0297.00043.
{\small\textit{Earliest verified question/agenda reference; historical priority qualified below. The introduction asks whether temporary jobs provide good careers or routes to permanent work, and then answers for 1990s Britain.}}

\item[\textnormal{[S2]}]\hypertarget{source:30007641:2}{} Mattia Filomena, Matteo Picchio. 2021. Are Temporary Jobs Stepping Stones or Dead Ends? A Meta-Analytical Review of the Literature. Section 1, Introduction, printed pp. 1–2; section 2, printed pp. 2–3.
\url{https://docs.iza.org/dp14367.pdf}.
DOI: 10.1108/IJM-02-2022-0064.
{\small\textit{Supporting or current literature; see evidence qualification. The review documents conflicting causal findings and context-dependent theoretical predictions, without a universal sign conclusion.}}
\end{itemize}
\end{document}
